\documentclass[11pt]{article}
\usepackage[margin=1in]{geometry}
\usepackage{amsmath, amssymb, amsthm}
\usepackage{hyperref}

\newtheorem{theorem}{Theorem}
\newtheorem{definition}{Definition}
\newtheorem{lemma}{Lemma}

\title{Information Geometry: The Fisher Metric}
\author{math-foundations / concept 40}
\date{}

\begin{document}
\maketitle

\section{Setting}

Let $\{p_\theta : \theta \in \Theta\}$ be a smoothly parameterized family of
probability densities on a measurable space $\mathcal{X}$, with
$\Theta \subseteq \mathbb{R}^d$ open. We assume sufficient regularity to
interchange differentiation in $\theta$ with integration in $x$, and that
$\log p_\theta(x)$ is twice continuously differentiable in $\theta$ for almost
every $x$. Define the \emph{score} $s_\theta(x) := \nabla_\theta \log p_\theta(x)$.

\begin{definition}[Fisher information matrix]
The Fisher information matrix at $\theta$ is the $d \times d$ matrix
\[
F(\theta)
  \;=\; \mathbb{E}_{X \sim p_\theta}\!\left[s_\theta(X)\, s_\theta(X)^\top\right]
  \;=\; -\,\mathbb{E}_{X \sim p_\theta}\!\left[\nabla_\theta^{\,2} \log p_\theta(X)\right].
\]
The two expressions agree under the regularity assumptions above.
\end{definition}

\begin{definition}[Statistical manifold]
A \emph{statistical manifold} is a pair $(\Theta, F)$ where $\Theta$ is a
smooth manifold parameterizing $\{p_\theta\}$ and $F$ is the Fisher information
matrix viewed as a Riemannian metric tensor: for $u, v \in T_\theta \Theta$,
$\langle u, v \rangle_\theta := u^\top F(\theta) v$.
\end{definition}

\begin{definition}[Natural gradient]
Given a smooth loss $L : \Theta \to \mathbb{R}$, the \emph{natural gradient} is
\[
\widetilde{\nabla} L(\theta) \;=\; F(\theta)^{-1} \nabla L(\theta),
\]
defined wherever $F(\theta)$ is positive definite. It is the steepest-descent
direction under the Fisher metric.
\end{definition}

\section{The score has mean zero}

\begin{lemma}\label{lem:score-mean}
$\mathbb{E}_{X \sim p_\theta}[s_\theta(X)] = 0.$
\end{lemma}

\begin{proof}
Since $\int p_\theta(x)\, dx = 1$ for every $\theta$, differentiating in $\theta$
and exchanging derivative and integral,
\[
0 \;=\; \nabla_\theta \!\int p_\theta(x)\, dx
  \;=\; \int \nabla_\theta p_\theta(x)\, dx
  \;=\; \int \frac{\nabla_\theta p_\theta(x)}{p_\theta(x)}\, p_\theta(x)\, dx
  \;=\; \mathbb{E}_{p_\theta}[s_\theta(X)]. \qedhere
\]
\end{proof}

\section{Fisher is the local KL metric}

\begin{theorem}\label{thm:fisher-kl}
The second-order Taylor expansion of
$\epsilon \mapsto D_{\mathrm{KL}}(p_\theta \,\|\, p_{\theta + \epsilon})$ near
$\epsilon = 0$ is
\[
D_{\mathrm{KL}}\!\left(p_\theta \,\|\, p_{\theta + \epsilon}\right)
   \;=\; \tfrac{1}{2}\, \epsilon^\top F(\theta)\, \epsilon
   \;+\; O\!\left(\|\epsilon\|^3\right).
\]
\end{theorem}

\begin{proof}
Define $\Phi(\epsilon) := D_{\mathrm{KL}}(p_\theta \,\|\, p_{\theta + \epsilon})
  = \mathbb{E}_{X \sim p_\theta}\!\left[\log p_\theta(X) - \log p_{\theta+\epsilon}(X)\right].$
Set $\ell(\epsilon, x) := \log p_{\theta + \epsilon}(x)$, so
$\Phi(\epsilon) = \mathbb{E}_{p_\theta}[\ell(0, X) - \ell(\epsilon, X)]$.

Taylor-expand $\ell(\epsilon, x)$ in $\epsilon$ at $0$ to second order, with
remainder $R(\epsilon, x) = O(\|\epsilon\|^3)$ uniformly on compact sets under
our smoothness hypotheses:
\[
\ell(\epsilon, x)
  \;=\; \ell(0, x)
  + \epsilon^\top \nabla_\epsilon \ell(0,x)
  + \tfrac{1}{2}\, \epsilon^\top \nabla_\epsilon^{\,2} \ell(0, x)\, \epsilon
  + R(\epsilon, x).
\]
Note $\nabla_\epsilon \ell(0,x) = s_\theta(x)$ and
$\nabla_\epsilon^{\,2} \ell(0,x) = \nabla_\theta^{\,2} \log p_\theta(x)$.
Substituting,
\[
\ell(0,x) - \ell(\epsilon, x)
  \;=\; -\,\epsilon^\top s_\theta(x)
  - \tfrac{1}{2}\, \epsilon^\top \nabla_\theta^{\,2} \log p_\theta(x)\, \epsilon
  - R(\epsilon, x).
\]
Take expectation under $p_\theta$. By Lemma~\ref{lem:score-mean},
$\mathbb{E}_{p_\theta}[\epsilon^\top s_\theta(X)] = \epsilon^\top \mathbb{E}_{p_\theta}[s_\theta(X)] = 0$,
so the linear term vanishes. By the Hessian form of Fisher,
$-\mathbb{E}_{p_\theta}[\nabla_\theta^{\,2} \log p_\theta(X)] = F(\theta)$, so
\[
\Phi(\epsilon)
  \;=\; 0 \;+\; \tfrac{1}{2}\, \epsilon^\top F(\theta)\, \epsilon
  \;-\; \mathbb{E}_{p_\theta}[R(\epsilon, X)]
  \;=\; \tfrac{1}{2}\, \epsilon^\top F(\theta)\, \epsilon \;+\; O\!\left(\|\epsilon\|^3\right). \qedhere
\]
\end{proof}

\paragraph{Consequence.} The Hessian of $\epsilon \mapsto
D_{\mathrm{KL}}(p_\theta \| p_{\theta+\epsilon})$ at $\epsilon = 0$ is
$F(\theta)$. The Fisher information is therefore the unique tensor that
captures the local quadratic behaviour of KL divergence — justifying its role
as the natural Riemannian metric on the statistical manifold and as the
preconditioner in natural-gradient descent.

\end{document}
